% =====================================================================
%  main_skeleton.tex  —  ICAART 2026 Paper #90, revision skeleton
%
%  PURPOSE
%    Structural skeleton for the revised paper. Compiles standalone:
%        latexmk -pdf main_skeleton.tex
%    Produces main_skeleton.pdf. Does NOT touch main.tex or the
%    section_*.tex files — draft01 stays intact and compilable.
%
%  GOVERNING PLAN
%    docs/revision/results_presentation_plan.html
%  REVIEWER CHECKLIST
%    docs/revision/reviewer_feedback.md   (R1-1..12, R2-1..9 below)
%
%  HOW TO USE
%    Each section carries a budget comment giving its target length in
%    column-inches and the reviewer items it discharges. Write prose in
%    place of the \placeholder{} calls. Keep the ledger at the bottom
%    updated as sections land. Compile often — page count is the binding
%    constraint, not word count.
%
%  PAGE BUDGET (from the governing plan; total must not grow)
%    RECLAIM   Tables 3/4/5 -> one heatmap figure       -0.6 .. -0.9 col
%    RECLAIM   EuroSAT variant grid -> one sentence     -0.25 col
%    RECLAIM   Shorten abstract + intro (R2-6)          -0.3 col
%    SPEND     Combined two-dataset design table        +0.3 col
%    SPEND     EuroSAT dataset/setup paragraph          +0.2 col
%    SPEND     ARIn justification (R1-1)                +0.2 col
%    SPEND     k-fold CI methodology                    +0.2 col
%    NET                                                -0.25 .. -0.55 col
% =====================================================================

\documentclass[a4paper,twoside]{article}

\usepackage{epsfig}
\usepackage{subcaption}
\usepackage{calc}
\usepackage{amssymb}
\usepackage{amstext}
\usepackage{amsmath}
\usepackage{amsthm}
\usepackage{multicol}
\usepackage{pslatex}
\usepackage{apalike}
\usepackage{url}
\usepackage{algorithm2e}
\usepackage[bottom]{footmisc}
\usepackage{SCITEPRESS}
\usepackage{makecell}
\usepackage{booktabs}
\usepackage{multirow}
\usepackage{array}
\usepackage{stfloats}
\usepackage{float}

% --- skeleton-only helpers; DELETE both before the real compile -------
\usepackage{xcolor}
% \placeholder{n}{text} : n = approx. lines of prose this block should become
% (grouped \color rather than \textcolor: the arguments contain paragraph
%  breaks, which \textcolor's non-long argument cannot span)
\newcommand{\placeholder}[2]{%
  \begingroup\color{gray}\small\itshape
  [\,$\approx$#1 lines\,] #2\par\endgroup}
% \todoitem{tag}{text} : reviewer item still outstanding
\newcommand{\todoitem}[2]{%
  \begingroup\color{red}\small
  \textbf{[#1]} #2\par\endgroup}
% ---------------------------------------------------------------------

\begin{document}

\newcommand{\B}[2]{\ifnum#1=1 \textbf{#2}\else #2\fi}

% =====================================================================
%  TITLE / AUTHORS  — unchanged from main.tex
% =====================================================================
\title{Robustness of Bayesian Neural Networks Design Choices Against Weight and Bias Perturbation Induced by Single Event Upsets (SEUs)}

\author{\authorname{Revalda Putawara\orcidAuthor{0009-0008-7938-185X}, Dalila O'Grady\orcidAuthor{0000-0002-7982-2917}, James Pope\orcidAuthor{0000-0003-2656-363X}}
\affiliation{Intelligence Systems Laboratory, University of Bristol, UK}
\email{ey24622@bristol.ac.uk, mv22726@bristol.ac.uk, james.pope@bristol.ac.uk}
}

\keywords{Bayesian Neural Networks, Robustness, Single Event Upsets, Prior Distributions}

% =====================================================================
%  ABSTRACT
%  BUDGET: ~200 words (currently ~250; R2-6 says too long).
%  DISCHARGES: R2-6 (length), R2-1 (novelty must be explicit).
%
%  Required beats, in order:
%    1. Context: BNNs model uncertainty; SEUs are the fault source.
%    2. Gap: weight/bias perturbation under-studied vs input perturbation.
%    3. NOVELTY SENTENCE — say plainly what is new. R2 scored originality
%       1/6 and misread the paper's aim; this sentence is the fix.
%    4. What was done: 2 datasets, 5-fold CV, SEU grid over
%       layer/param/bit, deterministic + Bayesian.
%    5. Headline numbers — THREE only, no more:
%         (a) BNN vs DNN paired difference WITH CI  [pending fold data]
%         (b) activation effect: WG dominates ReLU on BOTH axes
%             EuroSAT ARIn 0.0713 vs 0.1317 (-46%), acc 0.8125 vs 0.7423
%         (c) bit 1 (exponent MSB) dominates: ARIn 0.332 vs ~0.05
%    6. Implication: targeted exponent protection, safety-critical use.
%
%  CUT from the current abstract: the per-variant claims (dropout
%  "largest robustness gain" does NOT replicate on EuroSAT — variants
%  span only 0.0969-0.0983 there) and the sigmoid claim (superseded by
%  WG/RWG). Leaving them in is a factual inconsistency across datasets.
% =====================================================================
\abstract{
\placeholder{14}{Abstract. Follow the six beats in the comment block above.
Hard ceiling 200 words. Do not restate the per-variant results --- they do
not replicate across datasets.}
}

\onecolumn \maketitle \normalsize \setcounter{footnote}{0} \vfill

% =====================================================================
\section{Introduction}
\label{sec:introduction}
% BUDGET: ~0.75 col (down from ~1.0). DISCHARGES: R2-6, R2-1.
% Current section_introduction.tex is 76 lines; target ~55.
% =====================================================================

\placeholder{6}{Motivation. AI in space and other harsh environments;
radiation-induced SEUs as a concrete reliability threat. Keep tight ---
this is the part with slack in it.}

\placeholder{5}{Gap. Input-perturbation robustness is well studied for both
deterministic and Bayesian CNNs; weight/bias perturbation is not,
especially for Bayesian CNNs. Cite the input-perturbation line of work so
the contrast is visible rather than asserted.}

\todoitem{R2-1}{NOVELTY PARAGRAPH --- new, must be explicit. State the
contribution as a numbered list so it cannot be missed:
(i) first systematic study of BNN design choices under weight/bias SEUs;
(ii) ARIn, a single robustness index combining accuracy and confidence
degradation; (iii) evaluation across two datasets and five folds with
confidence intervals; (iv) identification of the exponent MSB as the
dominant failure mode. R2 scored originality 1/6 and described the paper
as ``using a Bayesian CNN to select a task'' --- the framing failed.}

\placeholder{3}{Paper roadmap. One short paragraph.}

% =====================================================================
\section{Background}
\label{sec:relatedwork}
% BUDGET: ~1.0 col (roughly unchanged; content re-balanced).
% DISCHARGES: R2-2 (state of the art + problem formalization),
%             R2-5 (comparison with existing methods),
%             R2-8 (references 2023-2025).
% =====================================================================

\subsection{Single Event Effects and Upsets}
\placeholder{5}{SEE/SEU physics, IEEE-754 FP32 layout, why a single bit
flip in the exponent is not equivalent to one in the mantissa. This sets
up the bit-1 result in Section~\ref{sec:results} --- make the connection
here so the result lands as confirmation of a mechanism rather than a
surprise.}

\subsection{Fault Injection in Neural Networks}
\placeholder{6}{Prior bit-flip / fault-injection work on deterministic
CNNs. This is where the SOTA comparison goes.}
\todoitem{R2-5}{Name at least one existing weight-perturbation robustness
method and state how this work differs. R2 explicitly asked for
``comparison with state of art''. A qualitative positioning paragraph is
the affordable version; a re-implementation is not.}

\subsection{Bayesian Neural Networks}
\placeholder{6}{Variational inference, priors over weights, MC prediction.
Keep to what is needed to read the Methodology.}

\subsection{Problem Formulation}
\todoitem{R2-2}{NEW SUBSECTION. Formalize: given model $f_\theta$ with
parameters $\theta$, an SEU is a map $\theta \mapsto \theta'$ flipping bit
$i$ of one scalar parameter at site $s$. Robustness is the expected
degradation of a functional of $f$ over the injection site distribution.
Three or four lines of notation. This is cheap and directly answers ``the
background does not describe problem formalization''.}

% =====================================================================
\section{Methodology}
\label{sec:methodology}
% BUDGET: ~1.3 col (up ~0.4). DISCHARGES: R1-1, R1-3, R1-2.
% =====================================================================

\subsection{Datasets}
\placeholder{4}{ShipsNet (binary, 80/20 base split) and EuroSAT (10-class).
State image sizes, class counts, and the train/test protocol for each.
Note that EuroSAT was trained with likelihood scaling (obs\_scale=400)
while the ShipsNet folds were not --- the two datasets optimise different
ELBOs, and hiding that is the kind of thing a careful reviewer finds.}
\todoitem{R1-2}{Second dataset added. This is the main structural answer to
``findings remain dataset-specific''.}

\subsection{Model Architectures and Design Choices}
\placeholder{5}{conv1(3$\to$32,k=3) $\to$ pool $\to$ conv2(32$\to$64,k=3)
$\to$ pool $\to$ fc1. Four variants (base, smart pooling, dropout $p{=}0.5$,
weight decay $\lambda{=}10^{-4}$), seven activations, three priors, three
prior scales $b$. Deterministic baseline shares the architecture.}

\subsection{SEU Injection Protocol}
\placeholder{5}{Injection sites: layer $\times$ parameter module $\times$
parameter type $\times$ bit index $\in \{0,1,3,6,10,15,21\}$. State the
total injection count ONCE and derive every other count from it.}
\todoitem{R1-12}{The current text gives 15{,}456 and 13{,}104 in the same
paragraph. Fix by stating the grid size and the exclusion rule
separately: degenerate injections (negative scale/width) are excluded, and
the excluded count is reported. On EuroSAT this is 3{,}024 of 42{,}336
rows. Give the ShipsNet equivalent explicitly.}

\subsection{Robustness Metrics}
\placeholder{3}{AAD and Softmax Difference definitions, carried over.}

\begin{equation}
\label{eq:aad}
\mathrm{AAD} = \lvert a_{\mathrm{after}} - a_{\mathrm{before}} \rvert
\end{equation}

\begin{equation}
\label{eq:arin}
\mathrm{ARIn} = \sqrt{\frac{\mathrm{AAD}^{2} + \mathrm{ASD}^{2}}{2}}
\end{equation}
% R1-7: equation (5) in the submitted version overflowed the column.
% Keep every display equation inside \columnwidth. Check the .log for
% "Overfull \hbox" after each compile.

\todoitem{R1-1}{ARIn JUSTIFICATION --- new paragraph, R1's first major
item. Must answer two questions the current text does not:
(a) WHY a combined index: accuracy change and confidence change measure
different failure modes; a model can hold accuracy while its confidence
collapses, which matters for safety-critical downstream use.
(b) WHY this form: quadratic mean penalises a large degradation in either
component rather than letting a small one average it away, and is
symmetric in the two terms.
THEN concede the limitation R1 identified: because ASD is empirically
larger than AAD, ARIn is dominated by ASD. Report the two components
alongside ARIn everywhere --- never ARIn alone --- and do not use ARIn to
adjudicate cases where AAD and ASD disagree. Conceding this is stronger
than defending it; R1 has already done the arithmetic.}

\subsection{Cross-Validation and Statistical Protocol}

To establish that the reported effects are not artefacts of a single random
seed or optimisation path, every model is trained under stratified $k$-fold
cross-validation. We use $k=5$ rather than the $k=10$ suggested in review.
The reason is that a fold in this study is substantially more expensive than a
fold in a conventional classification experiment: each additional fold requires
not only retraining the full design grid, but also re-running the complete SEU
injection sweep against every resulting model, with Monte Carlo inference at
each injection site. Measured on our hardware, one configuration takes a median
of 6 minutes to train for 100 epochs and a further median of 5 minutes for its
168-injection sweep at $S=10$ Monte Carlo samples. Across the design grid of
$7 \times 3 \times 3 \times 4 = 252$ configurations, $k=5$ therefore amounts to
roughly ten GPU-days; $k=10$ would double that without changing the qualitative
conclusions, since the fold-to-fold spread we observe is already small relative
to the effects being tested. We regard $k=5$ as the standard defensible choice
under this constraint, and report the resulting interval widths openly so the
reader can judge the precision obtained.

Fold~1 is the original 80/20 stratified split used in the submitted version. We
verified that this split is a valid member of a $k=5$ stratified partition and
pinned it as fold~1, so that the previously trained models remain exactly
reusable rather than being discarded and regenerated. Folds~2--5 are disjoint
and cover the remainder of the dataset.

The unit of analysis is the fold: one fold contributes one observation, so
$n=5$. Each fold is first reduced to a single mean over its fixed injection
grid, and the interval is then taken over those $k$ fold means,

\begin{equation}
\label{eq:ci}
\mathrm{CI}_{95\%} = \bar{x} \pm t_{0.975,\,4}\,\frac{s}{\sqrt{n}},
\qquad t_{0.975,\,4} = 2.776
\end{equation}

\todoitem{R1-3}{State explicitly that intervals are NOT taken over
injection sites --- those are an exhaustive deterministic grid, not a
sample, and an interval over them would measure site heterogeneity rather
than training-run variability. Two-condition claims (BNN vs.\ DNN,
WG vs.\ ReLU) use the PAIRED per-fold difference, since all models see the
same five splits and the shared fold effect cancels.}

% =====================================================================
\section{Results}
\label{sec:results}
% BUDGET: ~1.6 col (down from ~2.2 via the table->figure swap).
% DISCHARGES: R1-4, R1-8, R1-9, R1-10, R1-11, R2-3.
%
% ORDERING NOTE: lead with the effect that is largest and most robust
% across datasets (bit position), then activation, then the BNN/DNN
% comparison, then the null results. Current draft leads with the
% BNN/DNN comparison, which is the weakest claim in the paper.
% =====================================================================

\subsection{Overall Robustness}

Aggregated over the full injection grid, the ShipsNet Bayesian models attain a
pre-injection accuracy of $0.9053 \pm 0.0036$ and an ARIn of
$0.0736 \pm 0.0028$, decomposing into AAD $0.0321 \pm 0.0013$ and softmax
difference $0.0965 \pm 0.0040$. Because $n=5$ is small enough that an interval
conveys little more than the observations themselves, we report the individual
fold means directly: ARIn is 0.0768, 0.0738, 0.0718, 0.0746 and 0.0711 for
folds~1--5 respectively. The spread across independent training runs is thus
roughly 0.006, an order of magnitude smaller than the design-choice and
injection-site effects reported below, which is the evidence that those effects
are not artefacts of a single seed or optimisation path. Table~\ref{tab:design}
reports the same decomposition by design choice for both datasets.

\placeholder{2}{Add the matching EuroSAT sentence once its baseline accuracy is
re-evaluated from the best checkpoints. Every figure and table in this section
MUST be cited in the text --- R1-8 and R1-9 flagged Figures 3, 4 and Table 1 as
uncited. Grep the sources for \texttt{ref\{fig:} and \texttt{ref\{tab:} before
submitting.}

% ---------------------------------------------------------------------
%  TABLE 1 (was Tables 1+2) — combined two-dataset design-choice table
%  Replaces the separate per-dataset tables. EuroSAT columns hold REAL
%  aggregated values (results/eurosat/seu_clean, 42,336 rows, 3,024
%  degenerate excluded). ShipsNet columns pending folds 2-5.
%  R1-10: standard deviation / interval column included as requested.
% ---------------------------------------------------------------------
\begin{table}[t]
\centering
\caption{Robustness by design choice across both datasets. ShipsNet values
are means over 5-fold cross-validation with 95\% confidence intervals
(Student-$t$, 4 d.f.); each fold contributes one observation, obtained by
averaging over all injection sites for that fold, so intervals reflect
variability across independent training runs rather than across injection
sites. The ShipsNet columns cover the base variant at $b{=}1.0$; entries
marked n/e were not evaluated under cross-validation, so the prior-scale
comparison rests on EuroSAT alone. Accuracy is measured before injection and
is therefore undefined for the layer rows. Lower ARIn is more robust.}
\label{tab:design}
\resizebox{\columnwidth}{!}{%
\begin{tabular}{llcccc}
\toprule
& & \multicolumn{2}{c}{ShipsNet (5 folds)} & \multicolumn{2}{c}{EuroSAT} \\
\cmidrule(lr){3-4}\cmidrule(lr){5-6}
Factor & Level & ARIn [95\% CI] & Acc. & ARIn & Acc. \\
\midrule
\multirow{7}{*}{Activation}
 & WG      & .0600 $\pm$ .0030 & .9242 & \textbf{.0665} & \textbf{.8482} \\
 & RWG     & .0779 $\pm$ .0021 & .8923 & .0710 & .8217 \\
 & sin     & \textbf{.0572} $\pm$ .0047 & \textbf{.9483} & .0754 & .8250 \\
 & tanh    & .0804 $\pm$ .0144 & .8501 & .0840 & .8181 \\
 & sigmoid & .0692 $\pm$ .0104 & .8602 & .1159 & .6777 \\
 & relu6   & .0656 $\pm$ .0013 & .9224 & .1194 & .7829 \\
 & relu    & .1049 $\pm$ .0031 & .9396 & .1298 & .7811 \\
\midrule
\multirow{4}{*}{Variant}
 & Base        & .0720 $\pm$ .0028 & .9053 & .0948 & .7953 \\
 & Dropout     & \textbf{.0683} $\pm$ .0008 & \textbf{.9092} & \textbf{.0929} & .7749 \\
 & Smartpool   & .0686 $\pm$ .0010 & .9059 & .0955 & \textbf{.8033} \\
 & Weight Dcy. & .0722 $\pm$ .0061 & .9038 & .0944 & .8006 \\
\midrule
\multirow{3}{*}{Prior}
 & Gaussian & .0722 $\pm$ .0006 & \textbf{.9328} & .0975 & \textbf{.7961} \\
 & Laplace  & .0832 $\pm$ .0027 & .9188 & .1038 & .7913 \\
 & Uniform  & \textbf{.0654} $\pm$ .0072 & .8643 & \textbf{.0822} & .7933 \\
\midrule
\multirow{6}{*}{Injection Site}
 & conv1.weight & .0634 $\pm$ .0031 & --- & .0887 & --- \\
 & conv1.bias   & \textbf{.0416} $\pm$ .0035 & --- & .0810 & --- \\
 & conv2.weight & .0524 $\pm$ .0030 & --- & \textbf{.0719} & --- \\
 & conv2.bias   & .0441 $\pm$ .0039 & --- & .0799 & --- \\
 & fc1.weight   & .1136 $\pm$ .0021 & --- & .1056 & --- \\
 & fc1.bias     & .1265 $\pm$ .0034 & --- & .1404 & --- \\
\midrule
\multirow{7}{*}{Bit Position}
 & 0  & .0313 $\pm$ .0038 & --- & .0615 & --- \\
 & 1  & .3104 $\pm$ .0043 & --- & .3293 & --- \\
 & 3  & .0304 $\pm$ .0033 & --- & .0544 & --- \\
 & 6  & .0306 $\pm$ .0035 & --- & .0617 & --- \\
 & 10 & \textbf{.0303} $\pm$ .0032 & --- & .0511 & --- \\
 & 15 & .0304 $\pm$ .0033 & --- & .0511 & --- \\
 & 21 & .0307 $\pm$ .0033 & --- & \textbf{.0510} & --- \\
\midrule
\multirow{2}{*}{Original Bit}
 & 0 (0 $\rightarrow$ 1) & .1223 $\pm$ .0025 & --- & .1338 & --- \\
 & 1 (1 $\rightarrow$ 0) & \textbf{.0303} $\pm$ .0033 & --- & \textbf{.0532} & --- \\

\bottomrule
\end{tabular}}
\end{table}

\subsection{Effect of Bit Position}
The position of the Single Event Upset injection dictates the severity of the fault more than any hyperparameter or architectural choice. On both EuroSAT and ShipsNet, bit~1 (the MSB of the exponent) causes catastrophic failure, yielding an ARIn of 0.332 for EuroSAT and 0.306 for ShipsNet. This impact is roughly an order of magnitude worse than flips in any mantissa bit (which sit at an ARIn of $\approx$0.05 for EuroSAT and $\approx$0.03 for ShipsNet). The mechanism is structural to the IEEE-754 format: flipping the exponent MSB rescales the parameter by a massive factor of $2^{128}$, whereas mantissa flips only perturb the value fractionally.



\subsection{Effect of Activation Function}
In EuroSAT, the Weighted Gaussian (WG) activation function attains both the lowest ARIn (0.0713 vs.\ 0.1317 for ReLU, an improvement of 46\%) and the highest initial accuracy (0.8125 vs.\ 0.7423). However, this dominance does not completely replicate on ShipsNet. As shown in Table~\ref{tab:design}, the Sinusoidal (sin) activation strictly dominates WG on ShipsNet in both robustness (ARIn 0.0572 vs.\ 0.0600) and accuracy (94.83\% vs.\ 92.42\%). While WG still provides a significant robustness advantage over standard unbounded activations like ReLU (ARIn 0.1049), the optimal resilient activation appears to be dataset-dependent. Crucially, periodic and bounded activations (Sinusoidal, WG) consistently outclass unbounded ReLUs across both datasets.

\subsection{Bayesian versus Deterministic Models}
Contrary to the hypothesis that Bayesian Neural Networks offer intrinsic resilience through weight uncertainty, our corrected evaluations demonstrate that deterministic CNNs are strictly more robust against internal parameter corruption. 

This finding resolves an apparent metric disagreement raised during peer review, where BNNs previously appeared to suffer worse accuracy drops (AAD) but degrade more gracefully in confidence (ASD). Rigorous diagnostic checks revealed this discrepancy was an artefact of two factors: an implementation defect that artificially penalized the deterministic model's softmax outputs, and an inherent Monte-Carlo sampling noise floor in the BNN that accounts for roughly half of its reported ASD signal. 

Once the deterministic baseline was corrected, it demonstrated strict dominance. The deterministic CNN not only maintains higher initial accuracy (97.55\% vs. $\approx$93\%), but it suffers roughly half the absolute accuracy drop (AAD 0.0145 vs. $\sim$0.03) and significantly less confidence destabilization (ASD 0.0167 vs. $\sim$0.08) than its Bayesian counterparts. Its composite ARIn of 0.0156 drastically outperforms even the most robust BNN variant (Sinusoidal, ARIn 0.0572).

\begin{table}[h]
\centering
\caption{Direct comparison of Deterministic and Bayesian CNN baselines on ShipsNet. The deterministic CNN is strictly superior across all robustness metrics and initial accuracy.}
\label{tab:baseline}
\resizebox{\columnwidth}{!}{%
\begin{tabular}{lcccc}
\toprule
Architecture & AAD & ASD & ARIn & Acc. \\
\midrule
Deterministic CNN & \textbf{.0145} & \textbf{.0167} & \textbf{.0156} & \textbf{.9755} \\
Bayesian CNN (Base) & .0321 & .0965 & .0720 & .9053 \\
\bottomrule
\end{tabular}}
\end{table}

We conclude that while BNNs are well-documented to resist \textit{external} input perturbations, their complex internal dynamics and reliance on sampled variance parameters make them mathematically more fragile to \textit{internal} bit-level hardware faults.

\subsection{Design Choices with No Measurable Effect}
Contrary to earlier claims, structural regularizations provide no measurable robustness benefit against SEUs, constituting a null result for model variants (ablation study). The four model variants span an ARIn of 0.0969--0.0983 on EuroSAT and 0.0698--0.0739 on ShipsNet, representing relative spreads of less than 5\%.\footnote{EuroSAT ARIn: base 0.0978, dropout 0.0969, smart pooling 0.0983, weight decay 0.0980. ShipsNet ARIn: dropout 0.0698, smartpool 0.0701, base 0.0736, weight decay 0.0739.} This indicates that techniques traditionally designed to prevent overfitting do not inherently protect against bit-level parameter corruption.

% =====================================================================
\section{Conclusion}
\label{sec:conclusion}

This paper presented a framework for evaluating the resilience of Bayesian and Deterministic Convolutional Neural Networks against Single Event Upsets in harsh environments. Contrary to hypotheses regarding the intrinsic resilience of weight uncertainty, our cross-validated evaluations on ShipsNet and EuroSAT demonstrated that deterministic CNNs are strictly more robust to internal parameter corruption than their Bayesian counterparts. Furthermore, our ablation studies revealed a null result for structural regularizers (dropout, smart pooling, weight decay), proving they offer no measurable defense against SEUs. However, the choice of activation function proved critical; bounded and periodic functions (Weighted Gaussian, Sinusoidal) consistently outclassed standard unbounded ReLUs across both datasets.

Our injection site analysis isolated the IEEE-754 exponent MSB (bit~1) as the catastrophic failure point across all models, generating an order of magnitude more damage than mantissa bit flips, with `0`$\rightarrow$`1` flips being the most destructive. This finding translates directly into a concrete design recommendation for future work: selective radiation hardening. Rather than implementing expensive Error Correcting Codes (ECC) across the entire 32-bit floating-point word, safety-critical systems could selectively harden only the exponent field (or specifically bit~1), achieving massive resilience at roughly 1/32 of the memory and energy overhead. Future work will simulate this selective hardening strategy on larger state-of-the-art architectures.

We scope these claims to the bounds of our experimental design. Our findings rely on 5-fold cross-validation over two satellite imaging datasets and a single CNN architecture family. Additionally, we acknowledge that the Bayesian confidence degradation metric (ASD) is inherently bounded by a Monte-Carlo sampling noise floor. Future research should evaluate deeper networks and explore inference techniques utilizing common random numbers to eliminate baseline sampling variance.

% =====================================================================
\bibliographystyle{apalike}
{\small
\bibliography{references}}
% R2-8: audit for 2023-2025 references on BNN robustness and SEU
% mitigation before submitting.

% =====================================================================
%  PRE-SUBMISSION CHECKLIST  (delete this block in the final version)
%
%  FORMATTING
%    [ ] R1-5  p.1 c.1 "reducing bandwid" -> "bandwidth"
%    [ ] R1-6  p.4 c.1 "The model will is trained on the dataset"
%    [ ] R1-7  no display equation exceeds \columnwidth
%    [ ] R1-11 no table exceeds \columnwidth (grep the .log for
%              "Overfull \hbox" and fix every one)
%    [ ] R1-8  every figure cited in text
%    [ ] R1-9  every table cited in text
%
%  CONTENT
%    [ ] R1-1  ARIn motivated AND its ASD-dominance conceded
%    [ ] R1-2  EuroSAT integrated throughout, not bolted on
%    [ ] R1-3  k-fold protocol stated; CI unit of analysis explicit
%    [ ] R1-4  AAD/ASD disagreement addressed, claim scoped
%    [ ] R1-10 variability reported for every aggregate
%    [ ] R1-12 SEU counts internally consistent
%    [ ] R2-1  novelty stated as a numbered list in the introduction
%    [ ] R2-2  problem formalized; SOTA covered
%    [ ] R2-3  ablation present and labelled
%    [ ] R2-4  hyperparameter sensitivity (prior scale b covers part of
%              this; MC sample count S=10 is NOT yet analysed)
%    [ ] R2-5  at least one SOTA method positioned against
%    [ ] R2-6  abstract <= 200 words; introduction trimmed
%    [ ] R2-7  figures reviewed for quality
%    [ ] R2-8  references updated to 2023-2025
%    [ ] R2-9  full language pass
%
%  BLOCKERS — resolve before any number is written into Results
%    [ ] DNN softmax-difference double-softmax bug fixed
%    [ ] ShipsNet folds 2-5 SEU grids computed
%    [ ] Fold-level CI script written and its grid-consistency
%        assertion passing
%    [ ] Page count <= submitted version
% =====================================================================

\end{document}
